\begin{refsection}
%\section{Selección de contenidos}

%En este capítulo se describe la metodología empleada en la selección de contenidos y cursos en la nueva propuesta de estructura curricular.



%\section{Reuniones con expertos por áreas disciplinares}

\section{Pertinencia de un bloque común en el plan de estudios}

Por las razones expuestas en la sección anterior, y con la idea de tener una mayor flexibilidad curricular en la carrera de Matemáticas, se plantea un bloque común (core) al inicio de la carrera, que permita explorar y estudiar diferentes áreas de la matemática posteriormente. Según la discusión que se planteó en la Reunión de Departamento del 05 de mayo de 2021, la planificación de la secuencia de los contenidos y las formas en las cuales se pueden organizar debe plantearse en función de las cinco áreas curriculares establecidas en el perfil de egreso: habilidades operacionales, pensamiento abstracto, formalismo, pensamiento algorítmico y habilidades investigativas. En el desarrollo de las capacidades cognitivas del estudiante, se definió iniciar con un enfoque operacional que disminuye progresivamente, mientras que conforme se van adquiriendo destrezas operacionales se va introduciendo y aumentando el nivel de abstracción y formalismo esperado. En paralelo, se deben ir desarrollando el pensamiento algorítmico y habilidades investigativas de forma creciente según el año de estudio.\\

De esta forma, en el core se busca suplir diferentes áreas del conocimiento que deben ser parte de una formación sólida de cualquier persona que desea ejercer como profesional en Matemática. En este proceso de aprendizaje, Bolaños (2015) recomienda que la secuencia de contenidos se puede organizar mediante cursos, módulos, talleres, seminarios y combinaciones de estos (p.11). Esta nueva forma de trabajar en la carrera permitirá establecer puntos en común según las exigencias y necesidades del Departamento de Matemática.

En el perfil de egreso se estableció que

\begin{center}
\begin{minipage}{0.75\linewidth}
        \vspace{5pt}%margen superior de minipage
        {\small
                Los primeros ciclos de la carrera deberán contener una alta proporción de aprendizaje conductista, con pequeños matices de los enfoques cognitivo y crítico… En esta etapa temprana, el enfoque cognitivo puede irse implantando de manera paulatina mediante talleres y actividades lúdicas o de estudio independiente, incrustadas dentro de los cursos teórico-prácticos… Hacia la parte final de la carrera se busca que [la persona estudiante] realice cognitivamente los procesos de aprendizaje. (p.57-58).
        }
        \vspace{5pt}%margen inferior de la minipage
    \end{minipage}
\end{center}

Incluso, ya se había mencionado desde los espacios de reflexión previos que el estudiantado “lo primero que debe aprender no tiene que ver con los contenidos, es
necesario que aprenda formas de organizar el pensamiento en Matemática” (p. 60) (**Esta postura es la parte crucial para la comprensión de lo que se está haciendo en los cursos de primer año**). Esta es una visión que se comparte en el cuerpo docente y una forma de dar respuesta a esta pretensión es el establecimiento de un core en la etapa inicial de la formación en la carrera de Matemática.

Asimismo, después de la revisión de los programas de la carrera de Matemática en diferentes universidades del mundo (UNAM, MIT, UChicago, Paris-Saclay, por mencionar solo algunas), se evidenció que esta forma de proponer un plan de estudio es una práctica usual, por lo que se reflexionó acerca de la pertinencia para adoptar prácticas similares en nuestra universidad. 

Este bloque común debe permitir un desempeño favorable en diferentes énfasis o áreas de especialización. Esto va acorde con la renovación y actualización del cuerpo docente, y con la diversificación de intereses de investigación, lo que permite una amplia gama de posibilidades para el estudiantado.

\section{Áreas curriculares}

La presente sección tiene por finalidad dar una definición a las distintas áreas curriculares presentadas en la Propuesta para el rediseño del Bachillerato en Matemática. Antes que nada, conviene definir qué se entenderá por áreas curriculares. Según el Artículo 2 del Acuerdo 033 del Consejo Superior Universitario del 2007, las áreas curriculares son “conjunto de programas curriculares afines que pueden ser agrupados porque sus referentes epistemológicos pertenecen a un área común de conocimiento” (Artículo 2 del Acuerdo 033 del Consejo Superior Universitario, 2007). Por su parte, un programa curricular

\begin{center}
\begin{minipage}{0.8\linewidth}
        \vspace{5pt}%margen superior de minipage
        {\small
                Es un sistema abierto y dinámico compuesto por actividades, procesos, recursos, infraestructura, profesores, estudiantes, egresados, mecanismos de evaluación y estrategias de articulación con la sociedad, mediante el cual se desarrolla un proceso que busca cumplir ciertos objetivos de formación en los estudiantes a través de sus planes de estudio. El título académico es el reconocimiento que hace la sociedad, a través de la Universidad, del cumplimiento de dichos objetivos de formación por parte de un individuo. (Artículo 3 del Acuerdo 033 del Consejo Superior Universitario, 2007).
        }
        \vspace{5pt}%margen inferior de la minipage
    \end{minipage}
\end{center}


En esta sección se describen las áreas curriculares según el análisis realizado en el perfil de egreso y la reunión de departamento efectuada el 05 de mayo de 2021.

\begin{table}[H]
\begin{center}
\begin{tabular}{|lccc|}
\hline
\multicolumn{4}{|c|}{\cellcolor[HTML]{00C0F3} \textbf{Habilidades operacionales}} \\ \hline
\multicolumn{1}{|l|}{} & \multicolumn{1}{c|}{\cellcolor[HTML]{8ED8F8} \textbf{Año}} & \multicolumn{1}{c|}{\cellcolor[HTML]{8ED8F8} \textbf{Dominio}} & \cellcolor[HTML]{8ED8F8} \textbf{Contenidos} \\ \cline{2-4} 
\multicolumn{1}{|l|}{} & \multicolumn{1}{c|}{I} & \multicolumn{1}{c|}{Alto} & Alto \\ \cline{2-4} 
\multicolumn{1}{|l|}{} & \multicolumn{1}{c|}{II} & \multicolumn{1}{c|}{Alto} & Medio \\ \cline{2-4} 
\multicolumn{1}{|l|}{} & \multicolumn{1}{c|}{III} & \multicolumn{1}{c|}{Alto} & Bajo \\ \cline{2-4} 
\multicolumn{1}{|l|}{\multirow{-5}{*}{\begin{tabular}[c]{@{}l@{}}Conjunto de todas las destrezas que requiere la persona \\ estudiante para poder efectuar operaciones matemáticas\\ concretas de cierta temática (aritmética, funciones, teoría\\  de conjuntos, cálculo, probabilidad, etc,), en la que no \\ es necesario un pensamiento abstracto. \end{tabular}}} & \multicolumn{1}{c|}{IV} & \multicolumn{1}{c|}{Alto} & Bajo \\ \hline
\end{tabular}
\end{center}
\end{table}

\begin{table}[H]
\begin{center}
\begin{tabular}{|llll|}
\hline
\multicolumn{4}{|c|}{\cellcolor[HTML]{00C0F3} \textbf{Pensamiento Abstracto}} \\ \hline
\multicolumn{1}{|l|}{} & \multicolumn{1}{l|}{\cellcolor[HTML]{8ED8F8}} & \multicolumn{1}{l|}{\cellcolor[HTML]{8ED8F8}} & \cellcolor[HTML]{8ED8F8} \\
\multicolumn{1}{|l|}{} & \multicolumn{1}{l|}{\cellcolor[HTML]{8ED8F8}} & \multicolumn{1}{l|}{\cellcolor[HTML]{8ED8F8}} & \cellcolor[HTML]{8ED8F8} \\
\multicolumn{1}{|l|}{} & \multicolumn{1}{l|}{\multirow{-3}{*}{\cellcolor[HTML]{8ED8F8} \textbf{Año}}} & \multicolumn{1}{l|}{\multirow{-3}{*}{\cellcolor[HTML]{8ED8F8} \textbf{Dominio}}} & \multirow{-3}{*}{\cellcolor[HTML]{8ED8F8} \textbf{Contenidos}} \\ \cline{2-4} 
\multicolumn{1}{|l|}{} & \multicolumn{1}{l|}{I} & \multicolumn{1}{l|}{Medio} & Bajo \\ \cline{2-4} 
\multicolumn{1}{|l|}{} & \multicolumn{1}{l|}{II} & \multicolumn{1}{l|}{Medio} & Medio \\ \cline{2-4} 
\multicolumn{1}{|l|}{} & \multicolumn{1}{l|}{III} & \multicolumn{1}{l|}{Alto} & Alto \\ \cline{2-4} 
\multicolumn{1}{|l|}{\multirow{-7}{*}{\begin{tabular}[c]{@{}l@{}}- Nivel elevado del pensamiento en el cual convergen la\\ deducción, la síntesis, la interpretación y el análisis.\\ - Medio para la construcción del conocimiento teórico a\\ través del proceso de formación del concepto.\\ - Supone la capacidad de asumir un marco mental de \\ trabajo, en el que necesariamente se tengan claros muchos\\ conceptos para poder generar más conocimiento teórico.\end{tabular}}} & \multicolumn{1}{l|}{IV} & \multicolumn{1}{l|}{Alto} & Alto \\ \hline
\end{tabular}
\end{center}
\end{table}

\begin{table}[H]
\begin{center}
\begin{tabular}{|llll|}
\hline
\multicolumn{4}{|c|}{\cellcolor[HTML]{00C0F3} \textbf{Pensamiento Algorítmico}} \\ \hline
\multicolumn{1}{|l|}{} & \multicolumn{1}{l|}{\cellcolor[HTML]{8ED8F8} \textbf{Año}} & \multicolumn{1}{l|}{\cellcolor[HTML]{8ED8F8} \textbf{Dominio}} & \cellcolor[HTML]{8ED8F8} \textbf{Contenidos} \\ \cline{2-4} 
\multicolumn{1}{|l|}{} & \multicolumn{1}{l|}{I} & \multicolumn{1}{l|}{Bajo} & Bajo \\ \cline{2-4} 
\multicolumn{1}{|l|}{} & \multicolumn{1}{l|}{II} & \multicolumn{1}{l|}{Medio} & Medio \\ \cline{2-4} 
\multicolumn{1}{|l|}{} & \multicolumn{1}{l|}{III} & \multicolumn{1}{l|}{Alto} & Alto \\ \cline{2-4} 
\multicolumn{1}{|l|}{\multirow{-5}{*}{\begin{tabular}[c]{@{}l@{}}Capacidad de entender, ejecutar, evaluar y\\ crear algoritmos. El pensamiento que hace \\ uso de procesos recursivos.\end{tabular}}} & \multicolumn{1}{l|}{IV} & \multicolumn{1}{l|}{Alto} & Alto \\ \hline
\end{tabular}
\end{center}
\end{table}

\begin{table}[H]
\begin{center}
\begin{tabular}{|llll|}
\hline
\multicolumn{4}{|c|}{\cellcolor[HTML]{00C0F3} \textbf{Formalismo}} \\ \hline
\multicolumn{1}{|l|}{} & \multicolumn{1}{l|}{\cellcolor[HTML]{8ED8F8} \textbf{Año}} & \multicolumn{1}{l|}{\cellcolor[HTML]{8ED8F8} \textbf{Dominio}} & \cellcolor[HTML]{8ED8F8}\textbf{Contenidos} \\ \cline{2-4} 
\multicolumn{1}{|l|}{} & \multicolumn{1}{l|}{I} & \multicolumn{1}{l|}{Bajo} & Bajo \\ \cline{2-4} 
\multicolumn{1}{|l|}{} & \multicolumn{1}{l|}{II} & \multicolumn{1}{l|}{Medio} & Medio \\ \cline{2-4} 
\multicolumn{1}{|l|}{} & \multicolumn{1}{l|}{III} & \multicolumn{1}{l|}{Alto} & Medio \\ \cline{2-4} 
\multicolumn{1}{|l|}{\multirow{-5}{*}{\begin{tabular}[c]{@{}l@{}}Implica que el estudiantado pueda estructurar de tal manera el\\  pensamiento para poderlo comunicar de una forma que se \\ ajuste a las reglas de legitimación matemáticas establecidas. \\ Si bien es cierto, el formalismo no es lo que mueve a un \\ matemático, si que es algo de lo que no se puede prescindir.\end{tabular}}} & \multicolumn{1}{l|}{IV} & \multicolumn{1}{l|}{Alto} & Alto \\ \hline
\end{tabular}
\end{center}
\end{table}

\begin{table}[H]
\begin{center}
\begin{tabular}{|llll|}
\hline
\multicolumn{4}{|c|}{\cellcolor[HTML]{00C0F3} \textbf{Habilidades investigativas}} \\ \hline
\multicolumn{1}{|l|}{} & \multicolumn{1}{l|}{\cellcolor[HTML]{8ED8F8}} & \multicolumn{1}{l|}{\cellcolor[HTML]{8ED8F8}} & \cellcolor[HTML]{8ED8F8} \\
\multicolumn{1}{|l|}{} & \multicolumn{1}{l|}{\multirow{-2}{*}{\cellcolor[HTML]{8ED8F8} \textbf{Año}}} & \multicolumn{1}{l|}{\multirow{-2}{*}{\cellcolor[HTML]{8ED8F8} \textbf{Dominio}}} & \multirow{-2}{*}{\cellcolor[HTML]{8ED8F8} \textbf{Contenidos}} \\ \cline{2-4} 
\multicolumn{1}{|l|}{} & \multicolumn{1}{l|}{} & \multicolumn{1}{l|}{} &  \\
\multicolumn{1}{|l|}{} & \multicolumn{1}{l|}{\multirow{-2}{*}{I}} & \multicolumn{1}{l|}{\multirow{-2}{*}{-}} & \multirow{-2}{*}{Bajo} \\ \cline{2-4} 
\multicolumn{1}{|l|}{} & \multicolumn{1}{l|}{} & \multicolumn{1}{l|}{} &  \\
\multicolumn{1}{|l|}{} & \multicolumn{1}{l|}{\multirow{-2}{*}{II}} & \multicolumn{1}{l|}{\multirow{-2}{*}{-}} & \multirow{-2}{*}{Bajo} \\ \cline{2-4} 
\multicolumn{1}{|l|}{} & \multicolumn{1}{l|}{} & \multicolumn{1}{l|}{} &  \\
\multicolumn{1}{|l|}{} & \multicolumn{1}{l|}{\multirow{-2}{*}{III}} & \multicolumn{1}{l|}{\multirow{-2}{*}{Bajo}} & \multirow{-2}{*}{Bajo} \\ \cline{2-4} 
\multicolumn{1}{|l|}{} & \multicolumn{1}{l|}{} & \multicolumn{1}{l|}{} &  \\
\multicolumn{1}{|l|}{\multirow{-10}{*}{\begin{tabular}[c]{@{}l@{}}Conjunto de habilidades de diversa naturaleza, que empiezan \\ a desarrollarse desde antes de que el individuo tenga acceso a\\ procesos sistemáticos de formación para la investigación, que\\ en su mayoría no se desarrollan solo para posibilitar la \\ realización de las tareas propias de la investigación, pero que \\ han sido detectadas por formadores como habilidades cuyo \\ desarrollo. en el(la) investigador(a) en formación o en \\ funciones, es una contribución fundamental para potenciar \\ que pueda realizar investigación de buena calidad.\end{tabular}}} & \multicolumn{1}{l|}{\multirow{-2}{*}{IV}} & \multicolumn{1}{l|}{\multirow{-2}{*}{Medio}} & \multirow{-2}{*}{Medio} \\ \hline
\end{tabular}
\end{center}
\end{table}

\section{Lógica de selección de contenidos}

La selección de contenidos responde a una lógica de construcción que inició con la identificación del objeto de estudio, se moldea con los saberes, competencias y valores del perfil de egreso determinados a partir de la recopilación de información contextual de la carrera, y que se termina de pulir con el aporte de profesionales y docentes.

Con base en la recopilación de información que se hizo en los talleres y asignaciones de inicio del 2021, se reportan las subáreas de contenidos con una puntuación mayor o igual que 1 (en una escala de 0 a 3) con base en las respuestas obtenidas para los dos primeros años; ver Tablas 4, 5, 6 y 7. Según se indicó, la nota brindada indica el nivel que cada contenido aporta a la consecución del área curricular correspondiente.

\begin{table}[H]
\begin{center}
\caption{Puntajes obtenidos para el primer año con calificación mayor o igual a 1. Las entradas vacías son cero}
\label{Primeranno}
\end{center}
\begin{center}
\begin{tabular}{|l|c|c|c|c|}
\hline
\rowcolor[HTML]{00C0F3} 
\multicolumn{1}{|c|}{\cellcolor[HTML]{00C0F3} \textbf{Contenidos}} & \textbf{Oper} & \textbf{Abs} & \textbf{Form} & \textbf{Alg} \\ \hline
Demostraciones & 2.3 & 1.7 & 1.6 &  \\ \hline
Inducción & 2.7 & 1.1 & 1.3 & 1.4 \\ \hline
Lógica proposicional & 1.7 & 1.6 & 1.6 & 1.1 \\ \hline
Teoría de conjuntos & 1.8 & 1.5 & 1.6 &  \\ \hline
Combinatoria & 2.3 &  &  & 1.5 \\ \hline
Bases de programación & 1.5 &  &  & 1.3 \\ \hline
Precálculo & 2.5 & 1 & 1 & 1.3 \\ \hline
Cálculo & 2.2 &  &  & 1.3 \\ \hline
Álgebra lineal & 1.9 &  &  & 1.2 \\ \hline
Análisis complejo & 1.3 &  &  &  \\ \hline
Probabilidad & 1.4 &  &  &  \\ \hline
Estadística & 1.3 &  &  &  \\ \hline
Preliminares de álgebra & 2.3 & 1.5 & 1.2 & 1.5 \\ \hline
Álgebra lineal & 1.5 &  &  &  \\ \hline
Cuerpos & 1.0 &  &  &  \\ \hline
Divisibilidad factorización & 2.4 & 1.2 &  & 2.1 \\ \hline
Aritmética modular & 2.3 &  &  & 2.1 \\ \hline
Funciones aritméticas & 1.4 &  &  & 1.1 \\ \hline
Álgebras de Boole & 1.1 &  &  &  \\ \hline
Geometría Euclídea & 1.8 & 1.4 & 1.4 & 1.3 \\ \hline
Geometría analítica & 1.1 &  &  &  \\ \hline
\end{tabular}
\end{center}
\end{table}

Con respecto al primer año, se observa que los contenidos dominantes en el área de habilidades operaciones son temas del área de fundamentos (demostraciones, inducción, lógica proposicional, teoría de conjuntos, combinatoria), análisis (precálculo, cálculo), álgebra (álgebra lineal, preliminares de álgebra, cuerpos), teoría de números (divisibilidad, factorización, aritmética modular, funciones aritméticas) y geometría (geometría euclídea y analítica). Se plantea por tal razón tener un bloque de dos cursos introductorios que incluya estos temas, menos los contenidos de análisis que por los temas específicos se agruparan en dos cursos adicionales. En el área curricular del pensamiento algorítmico, dominan los contenidos de teoría de números, combinatoria, preliminares de álgebra, inducción, precálculo y cálculo, en los cuales se estudian algoritmos específicos que facilitan el desarrollo de dicha área

\begin{table}[H]
\begin{center}
\end{center}
\begin{center}
\begin{tabular}{|l|c|c|c|c|}
\hline
\rowcolor[HTML]{00C0F3} 
\multicolumn{1}{|c|}{\cellcolor[HTML]{00C0F3} \textbf{Contenidos}} & \textbf{Oper} & \textbf{Abs} & \textbf{Form} & \textbf{Alg} \\ \hline
Demostraciones & 1.5 & 2.1 & 2.4 &  \\ \hline
Inducción & 2 & 1.5 & 1.6 &  \\ \hline
Lógica proposicional & 1.2 & 1.4 & 1.5 &  \\ \hline
Teoría de conjuntos & 1.3 & 1.6 & 1.9 & 1 \\ \hline
Combinatoria & 1.6 & 1.5 & 1.1 & 1.4 \\ \hline
Bases de programación & 1.6 & 1.2 &  & 1.8 \\ \hline
Estructura de datos & 1.3 &  &  & 1.2 \\ \hline
Paradigmas computacionales & 1.1 &  &  & 1.1 \\ \hline
Bases de datos & 1.0 &  &  &  \\ \hline
Cálculo & 1.8 & 1.9 & 2.1 & 1.5 \\ \hline
Álgebra lineal & 2.1 & 2.4 & 2.2 & 1.9 \\ \hline
Espacios métricos & 1.3 & 1.1 & 1.1 &  \\ \hline
Ánalisis real & 1.3 & 1.2 & 1.2 &  \\ \hline
Análisis complejo & 1.3 &  &  &  \\ \hline
Probabilidad & 1.3 &  &  &  \\ \hline
Estadística & 1.4 &  &  & 1.1 \\ \hline
Análisis numérico & 1.3 &  &  & 1.4 \\ \hline
ODEs & 1.6 &  &  &  \\ \hline
Preliminares de álgebra & 1.3 & 1.5 & 1.7 & 1.3 \\ \hline
Álgebra lineal & 1.8 & 2.0 & 1.8 & 1.5 \\ \hline
Grupos & 1.2 & 1.1 & 1.2 &  \\ \hline
Anillos & \multicolumn{1}{l|}{1.1} & \multicolumn{1}{l|}{1.1} & \multicolumn{1}{l|}{1.2} & \multicolumn{1}{l|}{} \\ \hline
Cuerpos & \multicolumn{1}{l|}{1.1} & \multicolumn{1}{l|}{1.1} & \multicolumn{1}{l|}{1.2} & \multicolumn{1}{l|}{} \\ \hline
Espacios vectoriales & \multicolumn{1}{l|}{2.0} & \multicolumn{1}{l|}{2.3} & \multicolumn{1}{l|}{2.3} & \multicolumn{1}{l|}{1.1} \\ \hline
Geometría Euclídea & \multicolumn{1}{l|}{1.0} & \multicolumn{1}{l|}{1.2} & \multicolumn{1}{l|}{1.4} & \multicolumn{1}{l|}{} \\ \hline
Geometría analítica & \multicolumn{1}{l|}{1.4} & \multicolumn{1}{l|}{} & \multicolumn{1}{l|}{} & \multicolumn{1}{l|}{} \\ \hline
Topología (esp. Euclídeo) & \multicolumn{1}{l|}{1.3} & \multicolumn{1}{l|}{1.4} & \multicolumn{1}{l|}{1.2} & \multicolumn{1}{l|}{} \\ \hline
Topología (esp. métricos) & \multicolumn{1}{l|}{1.2} & \multicolumn{1}{l|}{1.1} & \multicolumn{1}{l|}{} & \multicolumn{1}{l|}{} \\ \hline
\end{tabular}
\end{center}
\caption{Puntajes obtenidos para el segundo año con calificación mayor o igual a 1. Las entradas vacías son cero}
\end{table}

En el caso del segundo año, en donde se espera un incremento en las áreas de pensamiento algorítmico, formalismo y pensamiento abstracto, se obtuvo una alta calificación en los temas de demostraciones, cálculo, álgebra lineal, espacios vectoriales, teoría de conjuntos e inducción.

\begin{table}[H]
\begin{center}
\end{center}
\begin{center}
\begin{tabular}{|l|c|c|c|c|}
\hline
\rowcolor[HTML]{00C0F3} 
\multicolumn{1}{|c|}{\cellcolor[HTML]{00C0F3} \textbf{Contenidos}} & \textbf{Oper}& \textbf{Abs} & \textbf{Form} & \textbf{Alg} \\ \hline
Demostraciones &  & 1.5 & 1.5 &  \\ \hline
Lógica proposicional &  &  &  & 1 \\ \hline
Combinatoria & 1.2 & 1 &  & 1.1 \\ \hline
Bases de programación & 1.6 & 1.2 &  & 1.8 \\ \hline
Estructura de datos & 1.5 &  &  & 1.3 \\ \hline
Paradigmas computacionales & 1 & 1 &  & 1 \\ \hline
Bases de datos & 1.1 &  &  & 1.1 \\ \hline
Aprendizaje automático & 1.2 &  &  & 1.3 \\ \hline
Espacios métricos & 1.1 & 2.3 & 1.9 &  \\ \hline
Ánalisis real & 1.2 & 2.4 & 2.0 &  \\ \hline
Análisis complejo & 1.5 & 2 & 1.8 &  \\ \hline
Probabilidad & 1.3 & 1.8 & 1.2 & 1.0 \\ \hline
Estadística & 1.5 & 1.2 & 1.3 & 1.3 \\ \hline
Análisis numérico & 1.9 & 1.3 & 1.3 & 2.3 \\ \hline
Análisis funcional &  & 1.6 & 1.3 &  \\ \hline
ODEs & 1.5 & 2.3 & 2.1 & 1.6 \\ \hline
PDEs & 1.6 & 1.1 & 1.1 & 1.3 \\ \hline
Modelación  matemática & 1.2 & 1.1 &  & 1.1 \\ \hline
Grupos & 1.8 & 2.8 & 2.1 & 1.6 \\ \hline
Anillos & 1.9 & 2.8 & 1.22.2 & 1.7 \\ \hline
Cuerpos & 1.7 & 2.7 & 2.1 & 1.4 \\ \hline
Álgebras & 1.3 & 2.1 & 1.8 &  \\ \hline
Módulos & 1.0 & 1.8 & 1.8 &  \\ \hline
Teoría de Galois & 1.9 & 2.6 & 2.2 & 2.1 \\ \hline
Teoría de códigos & 1.3 & 1.1 &  &  \\ \hline
Curvas elípticas &  & 1.1 &  &  \\ \hline
Álgebra de Boole &  & 1.0 & 1.0 &  \\ \hline
Cálculo de predicados & 1.6 & 1.5 & 1.5 & 1.0 \\ \hline
Teoría de Computabilidad & 1.1 & 1.2 & 1.1 & 1.2 \\ \hline
Teoría de Modelos & 1.4 & 1.3 & 1.4 &  \\ \hline
Teoría de Demostración & 1.0 & 1.2 & 1.2 &  \\ \hline
Geometría proyectiva &  & 1.5 & 1.4 &  \\ \hline
Geometría diferencial & 1.0 & 1.0 &  &  \\ \hline
Topología de espacios euclídeos & 1.4 & 1.7 & 1.5 &  \\ \hline
Topología de espacios métricos & 1.5 & 2.1 & 2.0 &  \\ \hline
Topología general & 1.4 & 2.0 & 1.8 &  \\ \hline
\end{tabular}
\end{center}
\caption{Puntajes obtenidos para el tercer año con calificación mayor o igual a 1. Las entradas vacías son cero.}
\end{table}

Para el tercer año, en pensamiento abstracto y formalismo dominan los contenidos de espacios métricos, análisis real y complejo, probabilidad, ecuaciones diferenciales ordinarias, grupos, anillos, cuerpos, álgebras, teoría de Galois, cálculo de predicados, geometría proyectiva y topología. En el pensamiento algorítmico destacan análisis numérico, programación, teoría de Galois y álgebra.

\begin{table}[H]
\begin{center}
\end{center}
\begin{center}
\begin{tabular}{|l|c|c|c|c|}
\hline
\rowcolor[HTML]{00C0F3} 
\multicolumn{1}{|c|}{\cellcolor[HTML]{00C0F3} \textbf{Contenidos}} & \textbf{Oper} & \textbf{Abs} & \textbf{Form} & \textbf{Alg} \\ \hline
Bases de programación & 1.0 &  &  & 1.1 \\ \hline
Aprendizaje automático &  &  &  & 1.0 \\ \hline
Espacios métricos &  & 1.6 & 1.4 &  \\ \hline
Análisis real & 1.1 & 2.3 & 2.0 & 1.0 \\ \hline
Análisis complejo & 1.1 & 1.6 & 1.5 &  \\ \hline
Probabilidad & 1.3 & 2.2 & 1.6 & 1.5 \\ \hline
Estadística & 1.0 & 1.5 & 1.3 & 1.1 \\ \hline
Análisis numérico & 1.4 & 1.2 & 1.5 & 2.0 \\ \hline
Análisis funcional &  & 2.1 & 1.7 &  \\ \hline
PDEs & 1.2 & 2.3 & 2.0 & 1.5 \\ \hline
Análisis armónico & 1.4 & 1.7 & 1.4 & 1.3 \\ \hline
Modelación matemática & 1.5 & 1.4 & 1.3 & 1.7 \\ \hline
Grupos &  & 1.4 & 1.4 &  \\ \hline
Anillos &  & 1.0 & 1.0 &  \\ \hline
Cuerpos &  & 1.0 & 1.0 &  \\ \hline
Álgebras & 1.2 & 1.9 & 1.6 & 1.0 \\ \hline
Módulos & 1.3 & 1.9 & 1.5 & 1.0 \\ \hline
Teoría de Galois &  & 1.0 &  &  \\ \hline
Álgebra homológica & 1.7 & 2.4 & 2.0 & 1.5 \\ \hline
Teoría de códigos &  & 1.1 &  &  \\ \hline
Teoría de autómatas &  & 1.1 &  &  \\ \hline
Criptografía de la llave pública & 1.0 &  &  & 1.0 \\ \hline
Curvas elípticas & 1.4 & 1.8 & 1.1 & 1.8 \\ \hline
Criptografía avanzada & 1.9 & 1.3 & 1.0 & 2.7 \\ \hline
Teoría analítica de números & 1.3 & 1.6 & 1.2 &  \\ \hline
Teoría de Computabilidad & 1.4 & 2.4 & 2.2 & 1.1 \\ \hline
Teoría de Modelos & 1.7 & 2.5 & 2.5 & 1.0 \\ \hline
Teoría de Demostración & 1.5 & 1.8 & 1.7 &  \\ \hline
Teoría descriptiva de conjuntos (ZF/ZFC) & 1.3 & 2.0 & 1.9 &  \\ \hline
Geometría proyectiva &  & 1.2 & 1.2 &  \\ \hline
Geometría diferencial & 1.7 & 2.1 & 1.8 &  \\ \hline
Geometría algebraica & 1.9 & 2.6 & 2.2 & 1.2 \\ \hline
Topología de espacios euclídeos &  & 1.0 &  &  \\ \hline
Topología de espacios métricos &  &  &  &  \\ \hline
Topología general &  & 1.6 & 1.2 &  \\ \hline
Homotopía & 1.7 & 2.8 & 2.3 & 1.1 \\ \hline
Homología & 1.6 & 2.7 & 2.3 &  \\ \hline
\end{tabular}
\end{center}
\caption{Puntajes obtenidos para el cuarto año con calificación mayor o igual a 1. Las entradas vacías son cero.}
\end{table}

Para el cuarto año, existe una mayor diversidad en temas sugeridos para el desarrollo de las áreas de pensamiento abstracto, algorítmico y formalismo. En general, destacan análisis real, probabilidad, análisis funcional, ecuaciones en derivadas parciales, álgebra homológica, teoría de computabilidad, teoría de modelos, teoría de conjuntos, geometría diferencial, geometría algebraica, homotopía y homología.

\section{Organización}

A continuación se describe la lógica en la organización de los cursos de la estructura curricular propuesta. De acuerdo con los criterios de organización y secuencia de contenidos, los cursos se agrupan por elementos en común ---las cinco áreas curriculares definidas en el perfil de egreso: habilidades operacionales, pensamiento abstracto, formalismo, pensamiento algorítmico y habilidades investigativas--- y se disponen en una secuencia vertical (por años y ciclos) y horizontal (dentro de un mismo ciclo lectivo), de modo que se garantice continuidad e integración entre las actividades formativas.

La organización sigue un diseño con tronco común en los primeros ciclos y, a partir del segundo año, el énfasis en matemática pura y el énfasis en matemática aplicada. En el desarrollo de las capacidades cognitivas se inicia con un enfoque operacional que disminuye progresivamente, mientras que el nivel de abstracción, formalismo y pensamiento algorítmico aumenta conforme avanza la carrera; las habilidades investigativas se refuerzan de forma transversal y se concentran con mayor intensidad en los últimos ciclos. Esta disposición gradual y armónica busca que el estudiantado logre una mejor comprensión y mayor profundización, conforme a lo establecido en la normativa universitaria y en los lineamientos de CONARE para la secuencia curricular.

\subsection{Cursos con énfasis operacional}

Con base en los niveles propuestos de las áreas curriculares, el primer año concentra el énfasis operacional, con dominio y contenidos altos en esta área. Se pretende así satisfacer las necesidades de los cursos posteriores y permitir una transición acorde con el perfil de ingreso y la madurez matemática del estudiantado. Si bien estos cursos tienen un alto componente operacional ---tanto en contenidos como en el dominio esperado---, también incorporan, de manera inicial, aspectos de pensamiento abstracto. La evaluación privilegia la habilidad operacional, y los contenidos sirven de introducción, en la medida de lo posible, al pensamiento abstracto y al formalismo.

En el ciclo I se ubican MA-0151 Fundamentos de álgebra, trigonometría y geometría analítica y MA-0152 Matemática Exploratoria. El primero desarrolla destrezas operacionales e intuición geométrica sobre álgebra, trigonometría y geometría vectorial y analítica, priorizando ejemplos concretos. El segundo inicia al estudiantado en el razonamiento inductivo y deductivo mediante resolución de problemas, con exploraciones en objetos discretos (teoría de números, combinatoria) y primeras nociones de pruebas sencillas de carácter operacional.

En el ciclo II, MA-0251 Introducción al Cálculo en una variable continúa el énfasis operacional e intuitivo sobre límites, derivadas e integrales, con trabajo centrado en ejemplos concretos y técnicas de cálculo. 

En el segundo año se mantiene un componente operacional alto en dominio, aunque los contenidos operacionales pasan a un nivel medio. MA-0361 Álgebra Lineal I enfatiza habilidades operacionales e intuición geométrica (sistemas, matrices, determinantes), e introduce progresivamente abstracción y formalismo mediante espacios vectoriales y transformaciones lineales. MA-0351 Cálculo en varias variables prioriza igualmente la práctica, la intuición y el pensamiento geométrico sobre sucesiones, series y cálculo diferencial e integral en varias variables. En cursos posteriores el componente operacional no desaparece, pero resulta secundario frente al énfasis en abstracción, formalismo, pensamiento algorítmico o investigación, según lo descrito en las áreas curriculares.

\subsection{Cursos con énfasis en pensamiento algorítmico}

El pensamiento algorítmico se incorpora de forma creciente: bajo en el primer año, medio en el segundo y alto a partir del tercero. Se contempla la inclusión de cursos relacionados con el uso de tecnologías y la implementación de algoritmos, en aras de una formación alineada con el perfil de egreso y con las prácticas emergentes (ciencia de datos, computación científica, modelación), con contenidos actualizados según los requerimientos de la matemática en distintas áreas.

En el tronco común, MA-0341 Matemática Computacional (tercer ciclo) orienta el trabajo a implementar pensamiento algorítmico mediante programación en un lenguaje de alto nivel, retomando contenidos de demostraciones, cálculo y álgebra lineal para ejecutar, analizar y explorar algoritmos, priorizando la intuición y la implementación sobre el formalismo. En el sexto ciclo, MA-0641 Análisis Numérico I profundiza en la resolución de problemas mediante métodos numéricos, con análisis de error, eficiencia, convergencia e implementación, y propicia habilidades investigativas a través de búsqueda bibliográfica e implementación de métodos aplicados.

En el énfasis en matemática aplicada, CA-0203 Herramientas de ciencia de datos I desarrolla destrezas algorítmicas y tecnológicas (hojas de cálculo, R, \LaTeX{}, Markdown, control de versiones), y MA-0647 Modelación Matemática integra algoritmos, simulación y validación de modelos. Asimismo, cursos del tronco o de los énfasis incorporan componentes algorítmicos puntuales ---por ejemplo, experimentación computacional en álgebra lineal o implementación de algoritmos en teoría de números--- sin que ese sea su eje principal. Las optativas de núcleo y temáticas en análisis numérico, computación científica, aprendizaje automático, optimización y análisis de datos permiten profundizar este énfasis en el cuarto año.

\subsection{Cursos con énfasis en pensamiento abstracto y formalismo}

En paralelo con lo anterior, en el primer año se incluye MA-0252 Introducción a las Demostraciones opera como puente entre la matemática predominantemente operacional y la más formal: retoma las intuiciones de exploración del primer ciclo y las articula con lógica, conjuntos, relaciones, funciones, teoría de números elemental e inducción, para desarrollar habilidades de demostración.

El pensamiento abstracto y el formalismo se introducen de forma paulatina y se consolidan a partir del segundo año (niveles medios), hasta alcanzar dominio y contenidos altos en el tercero y el cuarto. Los cursos de esta línea exigen estructurar el pensamiento conforme a las reglas de legitimación matemáticas, demostrar resultados y comunicar argumentos con rigor.

En el tronco común del segundo año, MA-0451 Principios de análisis I formaliza el análisis en una variable real (completitud, convergencia, diferenciación), asumiendo las intuiciones y destrezas operacionales de los cursos de cálculo y las bases de demostración del primer año. MA-0461 Álgebra Lineal II eleva el nivel de abstracción y formalismo (producto interno, teoría espectral, forma de Jordan), con demostraciones y resolución de problemas teóricos. En el tercer año, MA-0551 Principios de Análisis II continúa la formalización del análisis en varias variables; MA-0515 Ecuaciones Diferenciales aporta el sustento teórico de las ecuaciones ordinarias; y MA-0635 Topología General ofrece un tratamiento riguroso de espacios métricos y topológicos, continuidad, compacidad, conexidad y nociones de homotopía.

En el énfasis en matemática pura, MA-0496 Teoría de Números y MA-0471 Introducción a la Geometría Diferencial refuerzan demostración, abstracción y puente hacia estructuras algebraicas y geométricas más generales; MA-0561 Álgebra Abstracta I y MA-0661 Álgebra Abstracta II profundizan en anillos, módulos, grupos, cuerpos y teoría de Galois. En el cuarto año, cursos como MA-0702 Análisis Complejo y MA-0705 Análisis Real I consolidan el formalismo y la abstracción propios del análisis avanzado. En el énfasis en matemática aplicada, aunque el eje no es exclusivamente formalista, se exige dominio conceptual sólido (por ejemplo en probabilidad, estadística y modelación) que articula abstracción con transferencia a problemas concretos.

\subsection{Habilidades investigativas}

Según se externó por varios y varias docentes en la reunión de departamento efectuada el 05 de mayo de 2021, las habilidades investigativas deben estar presentes de forma transversal durante toda la carrera. A modo de ejemplo, en dicha reunión se sugirió asignar trabajos sobre temas relacionados con los cursos, búsquedas bibliográficas, proyectos, videos, búsquedas en internet, referenciar y citar, realizar bibliografías, asistencia a charlas y coloquios, exposiciones, entre otros. Esto no implica estudiar temas adicionales, sino abordar temas propios de los cursos mediante tareas que desarrollen las habilidades investigativas.

En la estructura propuesta, estas habilidades pasan de un nivel bajo en los primeros años a medio en el cuarto. De manera transversal, los programas de cursos del tronco y de los énfasis incorporan proyectos cortos, indagación bibliográfica, experimentación con software y comunicación oral o escrita de resultados. De forma específica, en los últimos ciclos se concentran actividades de formación orientadas a la práctica investigativa y profesional: MA-0780 Comunicación en las ciencias desarrolla escritura y expresión oral científica, ética y conducta responsable en investigación; y MA-0781 Práctica profesional en matemática pura o MA-0781 Práctica profesional en matemática aplicada permiten integrar conocimientos en un entorno académico (práctica de investigación) o profesional (aplicación en entes externos). En el énfasis en matemática aplicada, MA-0647 Modelación Matemática culmina con un proyecto de modelación que articula formulación, análisis, simulación y comunicación de resultados.

\subsection{Cursos optativos}

Los cursos optativos se organizan para ampliar la flexibilidad curricular respecto del plan vigente ---en el cual solo tres cursos optativos de matemática resultaban insuficientes frente al perfil de egreso--- y se agrupan en dos categorías: optativas de núcleo, con programas fijos en áreas representativas del objeto de estudio (análisis, álgebra, geometría, topología, probabilidad, estadística, lógica, optimización, sistemas dinámicos, entre otras), y optativas temáticas, que incluyen cursos de «Tópicos de\ldots» con contenidos abiertos por oferta según la especialidad docente y las necesidades formativas.

La oferta optativa es sustancialmente mayor que la del plan actual y se articula con el énfasis en matemática pura y el énfasis en matemática aplicada: en el primero predominan profundizaciones disciplinares teóricas; en el segundo se incluyen, además, opciones vinculadas a modelación, computación científica, análisis de datos, aprendizaje automático y, cuando corresponda, cursos de otras disciplinas afines que permitan trabajo interdisciplinario. De este modo, la organización de optativas refleja tanto la representatividad del campo disciplinar como la transferencia hacia contextos académicos y productivos contemplados en el perfil de egreso.

Los cursos optativos se dividen en dos tipos:
\begin{itemize}
    \item Optativos núcleo: Son aquellos cursos que, si bien son considerados optativos a un nivel de bachillerato o licenciatura, forman parte de temas que son ejes fundamentales en estudios más avanzados de matemática.
    \item Optativos temáticos: Se refiere a cursos que exploran temas que aunque sean importantes dentro de la disciplina no representan ejes fundamentales para estudios posteriores.
\end{itemize}

\printbibliography[heading=subbibliography,section=\the\value{refsection}]
\end{refsection}
